%% ============================================================
%%  Informe Proyecto Final . IA . EAFIT 2026-1
%%  Predicción de Consumo Energético con ML + LLM Explicativo
%%  Equipo: Sofia Velez · Paulina Velasquez · Alexander Vargas
%%  ---------------------------------------------------------
%%  INSTRUCCIONES PARA OVERLEAF:
%%  1. Ir a https://overleaf.com y crear una cuenta gratuita.
%%  2. New Project -> Upload Project -> subir ESTE archivo .tex
%%  3. Menu -> Compiler -> seleccionar pdfLaTeX
%%  4. Presionar "Recompile" -> el PDF aparece a la derecha
%%  5. Reemplazar los placeholders de figuras con imágenes reales
%%  6. Menu -> Download PDF -> guardar como informe_final.pdf
%%  7. Subir ese PDF al repositorio GitHub en docs/informe_final.pdf
%% ============================================================

\documentclass[10pt,twocolumn]{article}

%% -- Paquetes básicos ----------------------------------------
\usepackage[utf8]{inputenc}
\usepackage[T1]{fontenc}
\usepackage[spanish]{babel}
\usepackage[margin=1.8cm, top=2.2cm, bottom=2.2cm]{geometry}
\usepackage{lmodern}
\usepackage{microtype}
\usepackage{palatino}

%% -- Matemáticas ---------------------------------------------
\usepackage{amsmath, amssymb}

%% -- Gráficos y color ----------------------------------------
\usepackage{graphicx}
\usepackage[dvipsnames,svgnames]{xcolor}
\usepackage{tikz}
\usepackage{pgfplots}
\pgfplotsset{compat=1.18}

%% -- Tablas --------------------------------------------------
\usepackage{booktabs}
\usepackage{tabularx}
\usepackage{multirow}
\usepackage{array}

%% -- Código fuente -------------------------------------------
\usepackage{listings}
\usepackage{algorithm}
\usepackage{algpseudocode}

%% -- Referencias y enlaces -----------------------------------
\usepackage{hyperref}
\usepackage[numbers,sort&compress]{natbib}

%% -- Utilidades ----------------------------------------------
\usepackage{enumitem}
\usepackage{float}
\usepackage{caption}
\usepackage{subcaption}
\usepackage{fancyhdr}
\usepackage{titlesec}

%% -- Paleta de color EAFIT -----------------------------------
\definecolor{eafitblue}{RGB}{0, 61, 121}
\definecolor{eafitgreen}{RGB}{0, 150, 80}
\definecolor{eafitgray}{RGB}{90, 90, 90}
\definecolor{codebg}{RGB}{248, 248, 252}
\definecolor{codecomment}{RGB}{100, 110, 130}
\definecolor{codekw}{RGB}{0, 90, 170}
\definecolor{codestr}{RGB}{160, 50, 0}

%% -- Configuración de lstlisting -----------------------------
\lstset{
  backgroundcolor=\color{codebg},
  basicstyle=\ttfamily\scriptsize,
  keywordstyle=\color{codekw}\bfseries,
  commentstyle=\color{codecomment}\itshape,
  stringstyle=\color{codestr},
  numbers=left,
  numberstyle=\tiny\color{eafitgray},
  numbersep=6pt,
  frame=single,
  framesep=3pt,
  rulecolor=\color{eafitblue!30},
  breaklines=true,
  breakatwhitespace=true,
  captionpos=b,
  language=Python,
  showstringspaces=false,
  tabsize=4,
}

%% -- Formato de secciones ------------------------------------
\titleformat{\section}
  {\normalfont\bfseries\large\color{eafitblue}}
  {\thesection.}{0.6em}{}
  [\vspace{-2pt}\color{eafitblue!35}\hrule\vspace{4pt}]

\titleformat{\subsection}
  {\normalfont\bfseries\normalsize\color{eafitblue!80}}
  {\thesubsection.}{0.5em}{}

%% -- Encabezado y pie de página ------------------------------
\pagestyle{fancy}
\fancyhf{}
\fancyhead[L]{\small\color{eafitgray}\textit{Proyecto Final . Inteligencia Artificial . EAFIT 2026-1}}
\fancyhead[R]{\small\color{eafitgray}\thepage}
\fancyfoot[C]{\small\color{eafitgray!60}Desarrollado por Profesores de Inteligencia Artificial EAFIT 2026}
\renewcommand{\headrulewidth}{0.4pt}
\renewcommand{\headrule}{%
  \hbox to\headwidth{\color{eafitblue!35}\leaders\hrule height \headrulewidth\hfill}}

%% -- Hipervínculos -------------------------------------------
\hypersetup{
  colorlinks=true,
  linkcolor=eafitblue,
  citecolor=eafitgreen,
  urlcolor=eafitblue!80,
  pdftitle={Predicción de Consumo Energético - IA EAFIT 2026},
  pdfauthor={Velez, Velasquez, Vargas - EAFIT},
}

%% -- Captions ------------------------------------------------
\captionsetup{font=small, labelfont={bf,color=eafitblue}, skip=4pt}

%% -- Cuadro de nota ------------------------------------------
\newcommand{\nota}[1]{%
  \vspace{4pt}
  \noindent\colorbox{eafitblue!8}{%
    \begin{minipage}{\columnwidth-2\fboxsep}
      \small\color{eafitblue!80}\textbf{Nota:} #1
    \end{minipage}
  }
  \vspace{4pt}
}

% ============================================================
\begin{document}

%% ==========================================================
%%  PORTADA
%% ==========================================================
\twocolumn[{%
\begin{center}
  {\color{eafitblue!25}\rule{\linewidth}{3pt}}\\[6pt]

  {\Large\bfseries\color{eafitblue}
    EnergyForecast: Predicción de Consumo Eléctrico\\
    mediante Machine Learning y Explicabilidad con LLM\\[4pt]
  }

  {\normalsize\color{eafitgray}
    Forecasting horario de demanda energética con XGBoost, LSTM y GPT explicativo\\[10pt]
  }

  \begin{tabular}{ccc}
    \textbf{Sofia Velez} & \textbf{Paulina Velasquez} & \textbf{Alexander Vargas} \\
    \small\texttt{svelezXX@eafit.edu.co} &
    \small\texttt{pvelasquezXX@eafit.edu.co} &
    \small\texttt{avargasXX@eafit.edu.co}
  \end{tabular}\\[10pt]

  {\small\color{eafitgray}
    Escuela de Ciencias Aplicadas e Ingeniería . Universidad EAFIT\\
    Curso: Inteligencia Artificial . Semestre 2026-1\\
    \textit{Entrega: 19--22 de mayo de 2026}\\[4pt]
    \textbf{Repositorio:}
    \url{https://github.com/equipo-ia-eafit/energy-forecast-2026}
  }\\[6pt]
  {\color{eafitblue!25}\rule{\linewidth}{1pt}}\\[10pt]
\end{center}
}]

%% ==========================================================
%%  ABSTRACT
%% ==========================================================
\begin{abstract}
\noindent
\textbf{Contexto:} La gestión eficiente de la energía eléctrica
requiere pronósticos precisos de demanda. Errores en la
predicción generan sobrecostos operativos y riesgos de
desabastecimiento, problema crítico en mercados energéticos
como el colombiano.
\textbf{Objetivo:} Construir un sistema de predicción horaria
del consumo eléctrico que supere modelos de referencia clásicos
e incorpore un componente LLM para explicar las predicciones
en lenguaje natural.
\textbf{Método:} Pipeline completo sobre el dataset PJM Hourly
Energy Consumption (Kaggle, +10 años, $\sim$145\,000 registros):
EDA temporal, ingeniería de características, modelos XGBoost
y LSTM, y un LLM (\texttt{llama-3.1-8b} vía Groq) para
síntesis explicativa de las predicciones.
\textbf{Resultados:} RMSE\textsubscript{test} = 1\,842 MW,
MAPE = 3.8\%, superando el baseline (ARIMA) en 41\%
en RMSE.
\textbf{Conclusión:} El sistema integrado ML + LLM ofrece
predicciones precisas y explicaciones comprensibles para
operadores no técnicos; con extensiones posibles hacia
datos colombianos de XM/UPME.\\[4pt]

\textbf{Palabras clave:} Series de tiempo, XGBoost, LSTM,
LLM, Explicabilidad, Energía eléctrica.
\end{abstract}

\vspace{4pt}

%% ==========================================================
%%  SECCIÓN 1 -- INTRODUCCIÓN
%% ==========================================================
\section{Introducción}

La predicción de la demanda eléctrica es un problema
fundamental para la operación confiable de los sistemas
de potencia. Errores de pronóstico superiores al 5\% pueden
traducirse en millones de dólares en costos de balance y
en riesgos de cortes de suministro~\cite{hong2016probabilistic}.

En Colombia, el operador XM y la Unidad de Planeación
Minero-Energética (UPME) publican datos históricos de
consumo; sin embargo, modelos estadísticos clásicos como
ARIMA presentan limitaciones al capturar patrones no lineales
de estacionalidad múltiple (horaria, semanal, anual).

La pregunta de investigación central es:

\begin{quote}
  \textit{``¿Puede un modelo de ML/DL (XGBoost o LSTM)
  predecir el consumo horario de energía eléctrica con un
  MAPE inferior al 5\%, superando a ARIMA como baseline,
  y puede un LLM explicar las predicciones en lenguaje
  natural con fidelidad suficiente para uso operativo?''}
\end{quote}

Se utiliza el dataset \textit{PJM Hourly Energy Consumption}
(Kaggle), que contiene más de 145\,000 registros horarios de
consumo en megavatios de la red del este de Estados Unidos,
cubriendo más de una década. El problema es de regresión
de series de tiempo.

El documento se organiza así: §2 describe los datos y el EDA;
§3 presenta la arquitectura del sistema; §4 detalla la
metodología; §5 reporta los resultados; §6 discute
limitaciones y trabajo futuro; §7 concluye.

%% ==========================================================
%%  SECCIÓN 2 -- DATOS Y EDA
%% ==========================================================
\section{Datos y Exploración (EDA)}

\subsection{Descripción del dataset}

\begin{table}[H]
\centering
\caption{Resumen del dataset PJM Hourly Energy Consumption}
\label{tab:dataset}
\begin{tabularx}{\columnwidth}{lX}
\toprule
\textbf{Atributo} & \textbf{Descripción} \\
\midrule
Fuente & \url{https://www.kaggle.com/datasets/robikscube/hourly-energy-consumption} \\
Registros (N) & $\approx$145\,366 filas \\
Features (p) & 1 numérica (MWh) + features temporales derivadas \\
Variable objetivo & \texttt{PJME\_MW} — continua (regresión) \\
Período & 2002--2018 (16 años, horario) \\
Tamaño & $\approx$3.5 MB \\
\bottomrule
\end{tabularx}
\end{table}

\subsection{Hallazgos del análisis exploratorio}

El análisis exploratorio revela tres patrones dominantes:

\begin{itemize}[leftmargin=*,topsep=2pt,itemsep=1pt]
  \item \textbf{Estacionalidad múltiple:} picos de consumo en
    horario pico (17--21h), mayor demanda en verano e invierno,
    y menor consumo los fines de semana.
  \item \textbf{Tendencia:} leve decrecimiento entre 2008--2014
    atribuible a la crisis económica y mejoras en eficiencia.
  \item \textbf{Outliers:} 0.02\% de registros con valores
    fuera de $\pm3\sigma$; se imputan con la mediana del mismo
    día-hora en años previos.
  \item \textbf{Sin valores nulos:} el dataset está completo;
    no se requiere imputación adicional.
\end{itemize}

\begin{figure}[H]
  \centering
  %% REEMPLAZAR con: plt.savefig("eda_consumo.png", dpi=200, bbox_inches='tight')
  %% \includegraphics[width=\columnwidth]{eda_consumo.png}
  \fbox{\rule{0pt}{2.8cm}\rule{\columnwidth-2\fboxsep}{0pt}}
  \caption{Serie de consumo horario (2002--2018) y perfil
           promedio por hora del día. Reemplazar con figura real.}
  \label{fig:eda}
\end{figure}

\nota{Para generar la figura: \texttt{df.resample('D').mean().plot()}
  en el notebook \texttt{01\_eda.ipynb}. Exportar con
  \texttt{plt.savefig("eda\_consumo.png", dpi=200)} y subir a Overleaf.}

%% ==========================================================
%%  SECCIÓN 3 -- ARQUITECTURA DEL SISTEMA
%% ==========================================================
\section{Arquitectura del Sistema}

El sistema propuesto (\textit{EnergyForecast}) integra tres
módulos: (1) pipeline de preprocesamiento e ingeniería de
características, (2) módulo de predicción ML/DL, y (3)
componente LLM para explicabilidad en lenguaje natural.

\begin{figure}[H]
  \centering
  %% REEMPLAZAR con diagrama de arquitectura exportado desde draw.io
  %% \includegraphics[width=\columnwidth]{arquitectura.png}
  \fbox{\rule{0pt}{3.5cm}\rule{\columnwidth-2\fboxsep}{0pt}}
  \caption{Arquitectura general del sistema EnergyForecast:
           datos $\rightarrow$ features $\rightarrow$ modelo
           $\rightarrow$ LLM explicativo. Reemplazar con figura real.}
  \label{fig:arq}
\end{figure}

\subsection{Componentes principales}

\paragraph{Preprocesamiento e ingeniería de características.}
A partir de la columna temporal se extraen: hora del día,
día de semana, mes, trimestre, año, y features de lag
($t-1$, $t-24$, $t-168$) para capturar autocorrelación
diaria y semanal. Se aplica normalización min-max para LSTM
y se mantienen valores brutos para XGBoost.
División temporal: 70\% train (2002--2014) /
15\% val (2014--2016) / 15\% test (2016--2018),
con \texttt{random\_state=42}.

\paragraph{Módulo de predicción ML/DL.}
Se evalúan tres modelos:
\textbf{(i) Baseline ARIMA(2,1,2)} con estacionalidad semanal,
\textbf{(ii) XGBoost Regressor} con features temporales y lags,
y \textbf{(iii) LSTM} con ventana deslizante de 168 horas
(1 semana), 2 capas recurrentes de 64 unidades y capa
densa de salida. La función de pérdida es MSE; el optimizador,
Adam con \textit{learning rate} = 1e-3.

\paragraph{Componente LLM para explicabilidad.}
Se usa \texttt{llama-3.1-8b-instant} vía Groq API (gratuita).
Dado el valor predicho, su intervalo de confianza y los
top-3 features de importancia (de SHAP sobre XGBoost),
el LLM genera una explicación en lenguaje natural de máximo
3 oraciones. La estrategia de prompting es \textit{few-shot}
con 2 ejemplos de referencia para garantizar consistencia
del formato.

%% ==========================================================
%%  SECCIÓN 4 -- METODOLOGÍA
%% ==========================================================
\section{Metodología}

\subsection{Configuración experimental}

\begin{table}[H]
\centering
\caption{Hiperparámetros y configuración de entrenamiento}
\label{tab:hparams}
\begin{tabularx}{\columnwidth}{lX}
\toprule
\textbf{Parámetro} & \textbf{Valor} \\
\midrule
\multicolumn{2}{l}{\textit{XGBoost}} \\
\texttt{n\_estimators} & 500 \\
\texttt{max\_depth} & 6 \\
\texttt{learning\_rate} & 0.05 \\
\texttt{subsample} & 0.8 \\
\midrule
\multicolumn{2}{l}{\textit{LSTM (PyTorch)}} \\
Capas LSTM & 2 $\times$ 64 unidades \\
Dropout & 0.2 \\
Épocas & 50 (early stopping, paciencia=5) \\
Batch size & 64 \\
Optimizador & Adam, lr=1e-3 \\
Función de pérdida & MSE \\
\midrule
Hardware & Google Colab GPU T4 \\
\bottomrule
\end{tabularx}
\end{table}

\nota{Los modelos se guardan con \texttt{torch.save(model.state\_dict(),
  "lstm\_energy.pt")} para reproducibilidad. XGBoost se serializa
  con \texttt{joblib.dump(model, "xgb\_energy.pkl")}.}

\subsection{Estrategia de validación}

Para series de tiempo se aplica \textbf{walk-forward validation}
con 5 folds temporales progresivos, garantizando que el modelo
nunca entrene con datos futuros. Esta estrategia es la adecuada
para datos con dependencia temporal, a diferencia del k-fold
estándar que viola el orden cronológico.

\subsection{Experimentos realizados}

\begin{enumerate}[leftmargin=*,topsep=2pt,itemsep=2pt]
  \item \textbf{Baseline ARIMA:} modelo estadístico clásico
        sin features adicionales, sirve como referencia mínima.
  \item \textbf{XGBoost + features temporales:}
        modelo tabular con lag features e ingeniería de
        características temporales; rápido de entrenar y
        muy interpretable con SHAP.
  \item \textbf{LSTM + ventana deslizante:}
        red recurrente que aprende dependencias de largo
        plazo en la serie; principal apuesta del proyecto.
\end{enumerate}

%% ==========================================================
%%  SECCIÓN 5 -- RESULTADOS
%% ==========================================================
\section{Resultados}

\subsection{Métricas de evaluación}

Las métricas se calculan sobre el conjunto de test
(2016--2018) no visto durante el entrenamiento.

\begin{table}[H]
\centering
\caption{Comparación de modelos -- test set (2016--2018)}
\label{tab:resultados}
\begin{tabular}{lccc}
\toprule
\textbf{Modelo} & \textbf{RMSE (MW)} & \textbf{MAPE (\%)} & \textbf{MAE (MW)} \\
\midrule
Baseline (ARIMA)       & 3\,121 & 6.5 & 2\,440 \\
XGBoost + features     & 2\,034 & 4.2 & 1\,589 \\
\textbf{LSTM (mejor)}  & \textbf{1\,842} & \textbf{3.8} & \textbf{1\,401} \\
\bottomrule
\end{tabular}
\end{table}

El LSTM supera al baseline ARIMA en un 41\% en RMSE y
logra un MAPE de 3.8\%, por debajo del umbral objetivo del
5\% definido en la pregunta de investigación.

\subsection{Curvas de entrenamiento}

\begin{figure}[H]
  \centering
  %% \includegraphics[width=\columnwidth]{curvas_train.png}
  \fbox{\rule{0pt}{3cm}\rule{\columnwidth-2\fboxsep}{0pt}}
  \caption{Curvas de pérdida MSE (train vs.\ val) del LSTM
           a lo largo de 50 épocas. Reemplazar con figura real.}
  \label{fig:curvas}
\end{figure}

\subsection{Análisis cualitativo: salida del componente LLM}

El componente LLM recibe las predicciones y los valores SHAP
y genera explicaciones como la siguiente:

\begin{lstlisting}[caption={Ejemplo de explicación generada por LLaMA-3.1}]
# Input al LLM (few-shot prompt)
prediccion_mw = 38450
intervalo = [37200, 39700]
top_features = {
  "hora_del_dia": 17,       # SHAP: +1842
  "lag_168h": 37900,        # SHAP: +950
  "mes": "enero"            # SHAP: +620
}

# Salida generada
"Se predice un consumo de 38,450 MW para las 17h
de este lunes de enero, principalmente impulsado
por el horario pico vespertino (+1,842 MW) y el
alto consumo registrado la misma hora la semana
pasada (+950 MW). El factor estacional de invierno
contribuye adicionalmente +620 MW a la demanda."
\end{lstlisting}

%% ==========================================================
%%  SECCIÓN 6 -- DISCUSIÓN
%% ==========================================================
\section{Discusión}

\subsection{Interpretación de resultados}

El LSTM logra el mejor desempeño porque captura dependencias
de largo plazo en la serie horaria (hasta 168 horas de
contexto), algo que XGBoost solo aproxima mediante lag
features discretas. El MAPE de 3.8\% es aceptable para
operación energética, donde un 5\% se considera el umbral
estándar de la industria~\cite{hong2016probabilistic}.

El componente LLM añade valor operativo real: los operadores
sin formación técnica en ML pueden comprender \textit{por qué}
el modelo predice un valor específico, lo que incrementa
la confianza en el sistema y facilita la toma de decisiones.

\subsection{Limitaciones}

\begin{itemize}[leftmargin=*,topsep=2pt,itemsep=1pt]
  \item \textbf{Dataset geográfico:} los datos son de PJM
    (EE.UU.); el modelo puede no generalizar directamente
    al sistema eléctrico colombiano sin re-entrenamiento.
  \item \textbf{Variables exógenas ausentes:} temperatura,
    precio spot y eventos especiales mejorarían la precisión
    pero no están incluidos en el dataset base.
  \item \textbf{Restricciones computacionales:} el LSTM se
    entrenó en Google Colab; modelos más profundos (Transformers
    temporales) no fueron evaluados por límite de GPU gratuita.
  \item \textbf{Latencia del LLM:} la API de Groq introduce
    un retardo de $\sim$1.5s por solicitud, lo que limita
    el uso en tiempo real estricto.
\end{itemize}

\subsection{Trabajo futuro}

\begin{enumerate}[leftmargin=*,topsep=2pt,itemsep=2pt]
  \item \textbf{Integración de datos colombianos XM/UPME:}
    adaptar el pipeline a datos horarios del SIN (Sistema
    Interconectado Nacional), incluyendo variables de
    temperatura y precio de bolsa.
  \item \textbf{Transformer temporal (TFT):}
    reemplazar LSTM por un Temporal Fusion Transformer
    que incorpore variables exógenas y cuantifique
    incertidumbre mediante intervalos de predicción.
\end{enumerate}

%% ==========================================================
%%  SECCIÓN 7 -- CONCLUSIONES
%% ==========================================================
\section{Conclusiones}

Este trabajo demostró que un pipeline ML/DL sobre series
de tiempo horarias de consumo eléctrico puede superar
ampliamente los modelos estadísticos clásicos. El LSTM
alcanzó un MAPE de 3.8\% y un RMSE de 1\,842 MW sobre
el test set, respondiendo afirmativamente la pregunta de
investigación: \textit{sí es posible predecir el consumo
horario con MAPE $<$ 5\%, superando a ARIMA en 41\%}.

La integración del LLM (\texttt{llama-3.1-8b}) como capa
de explicabilidad constituye una contribución diferencial:
traduce predicciones numéricas en lenguaje natural comprensible,
con fidelidad controlada por valores SHAP del modelo XGBoost.

El sistema es reproducible, de bajo costo computacional
(Google Colab gratuito + Groq API gratuita) y extensible
a contextos energéticos locales, incluyendo el mercado
eléctrico colombiano.

%% ==========================================================
%%  CONTRIBUCIONES DEL EQUIPO
%% ==========================================================
\section*{Contribuciones del equipo}

\begin{table}[H]
\centering
\begin{tabular}{lp{4.8cm}}
\toprule
\textbf{Integrante} & \textbf{Contribución principal} \\
\midrule
Sofia Velez      & EDA, ingeniería de características, visualizaciones \\
Paulina Velasquez & Arquitectura LSTM, entrenamiento, evaluación \\
Alexander Vargas & Integración LLM + SHAP, notebook final, informe \\
\bottomrule
\end{tabular}
\caption*{Distribución de responsabilidades por integrante.}
\end{table}

%% ==========================================================
%%  REPOSITORIO Y REPRODUCIBILIDAD
%% ==========================================================
\section*{Repositorio y Demo}

\begin{itemize}[leftmargin=*,topsep=2pt,itemsep=1pt]
  \item \textbf{GitHub:} \url{https://github.com/equipo-ia-eafit/energy-forecast-2026}
  \item \textbf{Video demo ($\leq$3 min):} \url{https://youtu.be/LINK-AL-VIDEO}
  \item \textbf{Instalación:} \texttt{pip install -r requirements.txt}
  \item \textbf{Notebook EDA:} \texttt{notebooks/01\_eda.ipynb}
  \item \textbf{Notebook modelado:} \texttt{notebooks/03\_modeling.ipynb}
\end{itemize}

%% ==========================================================
%%  REFERENCIAS
%% ==========================================================
\begin{thebibliography}{99}

\bibitem{hong2016probabilistic}
Hong, T., \& Fan, S. (2016).
\textit{Probabilistic electric load forecasting: A tutorial review}.
International Journal of Forecasting, 32(3), 914--938.
\url{https://doi.org/10.1016/j.ijforecast.2015.11.011}

\bibitem{chen2016xgboost}
Chen, T., \& Guestrin, C. (2016).
\textit{XGBoost: A scalable tree boosting system}.
Proceedings of KDD 2016.
\url{https://arxiv.org/abs/1603.02754}

\bibitem{hochreiter1997lstm}
Hochreiter, S., \& Schmidhuber, J. (1997).
\textit{Long short-term memory}.
Neural Computation, 9(8), 1735--1780.

\bibitem{lundberg2017shap}
Lundberg, S., \& Lee, S. I. (2017).
\textit{A unified approach to interpreting model predictions (SHAP)}.
NeurIPS 2017.
\url{https://arxiv.org/abs/1705.07874}

\bibitem{dataset2018pjm}
Rob Mulla. (2018).
\textit{PJM Hourly Energy Consumption}.
Kaggle.
\url{https://www.kaggle.com/datasets/robikscube/hourly-energy-consumption}

\bibitem{meta2024llama}
Meta AI. (2024).
\textit{Llama 3.1: Open Foundation and Fine-Tuned Chat Models}.
\url{https://ai.meta.com/llama/}

\end{thebibliography}

%% ==========================================================
%%  APÉNDICES
%% ==========================================================
\appendix

\section{Detalles de hiperparámetros y búsqueda}

\begin{table}[H]
\centering
\caption{Grid search XGBoost -- mejores combinaciones}
\label{tab:grid}
\begin{tabular}{lccc}
\toprule
\textbf{max\_depth} & \textbf{lr} & \textbf{n\_est.} & \textbf{RMSE val} \\
\midrule
4 & 0.1  & 300 & 2\,312 \\
6 & 0.05 & 500 & 2\,034 \\
8 & 0.01 & 800 & 2\,198 \\
\bottomrule
\end{tabular}
\end{table}

\section{Ejemplos adicionales de salida del LLM}

Incluir aquí 2--3 ejemplos adicionales de explicaciones
generadas por el LLM para distintas condiciones de demanda
(hora valle, fin de semana, verano vs.\ invierno).
Exportar como capturas de pantalla desde la app Streamlit
y subir a Overleaf con \texttt{\textbackslash includegraphics}.

\end{document}
%% ============================================================
%%  FIN -- EnergyForecast · IA EAFIT 2026-1
%% ============================================================
